%% ============================================================
%%  CHAPITRE 2 — Conception de la couche d'agents
%% ============================================================
\chapter{Conception de la couche d'agents}
\label{chap:conception}

%% Listes d'outils \mcptool{...} très denses : ouvrir les marges
%% intérieures pour autoriser les coupures à _ , ( etc.
\sloppy

L'introduction a posé le problème (redondance, absence d'ancrage factuel,
absence de séparation des privilèges) et le chapitre~1 a établi les fondations
théoriques : orchestration d'agents, Model Context Protocol et \gls{rag}
agentique. Le présent chapitre traduit ces fondations en une conception logique
complète, puis relit les implémentations livrées comme un \emph{catalogue
comparé de formes d'agents}. Il prépare ainsi la stratégie d'orchestration qui
fera l'objet du chapitre~3.

La première partie traite de l'indexation des activités, pierre angulaire sans
laquelle la non-redondance ne serait qu'un v{\oe}u. La deuxième présente le
protocole \gls{mcp} tel qu'il s'applique au système, avec le catalogue d'outils
qu'il expose et le principe du moindre privilège qui gouverne leur
distribution. La troisième détaille le catalogue des agents retenus, en
réservant une matrice à quatre variantes pour l'agent de création de syllabus.
La quatrième propose une lecture comparative des implémentations : non un test
de conformité, mais un examen architectural de ce que chaque variante
privilégie et de ce qu'elle paie en contrepartie. La conclusion enchaîne sur le
chapitre~3, en montrant que ces différences de conception appellent une
stratégie d'orchestration de niveau supérieur.

%% ────────────────────────────────────────────────────────────
\section{Indexation des activités pédagogiques}
\label{sec:indexation}

\subsection{Le problème de la redondance}

L'approche mono-modèle décrite en introduction souffre d'un défaut structurel :
chaque appel au \gls{llm} produit du contenu nouveau sans mémoire de ce qui a
déjà été produit. Sur un parcours pédagogique enrichi au fil de plusieurs
sessions, plusieurs activités finissent par traiter le même concept sous des
angles à peine différents.

Cette redondance n'est pas qu'un problème esthétique. Elle dilue l'ancrage des
objectifs d'apprentissage, alourdit la charge cognitive de l'apprenant et
complique le travail de l'enseignant qui doit trier manuellement les doublons.
Pour qu'un agent génère du contenu utile, il doit d'abord savoir ce qui existe.

La solution retenue est un mécanisme d'indexation qui transforme chaque
activité existante en une représentation numérique comparable, et qui expose
une opération de recherche par similarité aux agents générateurs.
L'indexation constitue la pierre angulaire du système : sans elle, les parties
suivantes (\gls{mcp}, agents) perdraient leur raison d'être.

\subsection{Représentation vectorielle des activités}

Le principe est le suivant. Chaque activité pédagogique (cours, question,
exercice ou quiz) est décrite par un texte : son titre concaténé avec son
contenu. Ce texte est transformé en un vecteur numérique de dimension fixe
par un modèle de plongement (\emph{embedding model}). La dimension du vecteur
est déterminée par le modèle retenu ; elle est sans incidence sur la logique
de l'agent, qui ne manipule jamais les composantes individuelles. Deux textes
sémantiquement proches, par exemple deux exercices portant sur la résolution
d'équations du second degré, produisent des vecteurs proches dans cet espace,
même si leur formulation diffère.

\begin{figure}[H]
\centering
\begin{tikzpicture}[
    node/.style={rectangle, rounded corners=6pt, draw, thick,
      minimum width=3.0cm, minimum height=1.0cm, align=center},
    arrow/.style={->, thick, -{Stealth[length=3mm]}}
]
\node[node, fill=myblue!8] (txt) {Titre + contenu\\de l'activité};
\node[node, right=1.5cm of txt, fill=myblue!8] (emb) {Modèle de\\plongement};
\node[right=1.6cm of emb] (vec) {$\vec{v}\in\mathbb{R}^{d}$};
\draw[arrow] (txt) -- (emb);
\draw[arrow] (emb) -- (vec);
\end{tikzpicture}
\caption{Transformation d'une activité en vecteur numérique.}
\label{fig:embedding}
\end{figure}

À partir de cette représentation, le système définit une \textbf{mesure de
similarité} entre deux activités. Soit $\vec{u}$ le vecteur d'une activité
existante et $\vec{v}$ celui d'un brouillon à créer ; la similarité
$s(\vec{u}, \vec{v})$ est une fonction qui croît lorsque les deux vecteurs
pointent dans la même direction de l'espace de plongement, et qui décroît
lorsqu'ils en divergent. Le choix précis de cette mesure (cosinus, distance
euclidienne ou autre) relève de la configuration du moteur d'indexation et
non de la conception de l'agent.

Une calibration de cette similarité — un point au-delà duquel deux activités
sont considérées comme un quasi-doublon — peut être fixée côté serveur de
manière empirique. Cette calibration reste un paramètre d'implémentation~:
elle n'appartient pas au contrat de l'agent, qui consomme la primitive de
recherche sans en connaître les constantes. Le chapitre ne fait donc pas
reposer son argumentaire sur une valeur numérique particulière~; ce sont les
\emph{propriétés} de la recherche par similarité (ordonnancement par
proximité, signal complémentaire de la lecture directe) qui structurent la
conception.

\subsection{Pipeline d'indexation}
\label{sec:pipeline-index}

Le pipeline d'indexation comporte deux flux complémentaires : la
\textbf{création} d'un index à partir des activités existantes, et la
\textbf{recherche} par similarité avant chaque nouvelle écriture.

\begin{figure}[H]
\centering
\begin{tikzpicture}[
    node/.style={rectangle, rounded corners=5pt, draw, thick,
      minimum width=3.0cm, minimum height=0.8cm, align=center,
      font=\small},
    db/.style={cylinder, draw, thick, shape border rotate=90,
      aspect=0.3, minimum height=1.2cm, minimum width=1.8cm,
      align=center, font=\small, fill=myblue!8},
    arrow/.style={->, thick, -{Stealth[length=3mm]}}
]
%% Flux gauche : création de l'index
\node[node, fill=myblue!8] (act) {Activité\\(titre + contenu)};
\node[node, below=0.7cm of act, fill=myblue!8] (mod1) {Modèle de\\plongement};
\node[db, below=0.7cm of mod1] (db1) {Index\\vectoriel};
\draw[arrow] (act) -- (mod1);
\draw[arrow] (mod1) -- (db1);
\node[above=0.15cm of act, font=\bfseries\small] {Création};

%% Flux droite : recherche
\node[node, right=3.0cm of act, fill=myblue!8] (q) {Requête\\(sujet)};
\node[node, below=0.7cm of q, fill=myblue!8] (mod2) {Modèle de\\plongement};
\node[node, below=0.7cm of mod2, fill=myblue!8] (sim)
  {Recherche par\\similarité ($k$-NN)};
\node[node, below=0.7cm of sim, fill=myblue!8] (res)
  {Top-$k$ activités\\les plus proches};
\draw[arrow] (q) -- (mod2);
\draw[arrow] (mod2) -- (sim);
\draw[arrow] (sim) -- (res);
\draw[arrow] (db1) -- node[above, font=\footnotesize]{interroge} (sim);
\node[above=0.15cm of q, font=\bfseries\small] {Recherche};
\end{tikzpicture}
\caption{Pipeline d'indexation : création (gauche) et recherche (droite).}
\label{fig:pipeline}
\end{figure}

À la création, chaque nouvelle activité passe par le modèle de plongement et
le vecteur résultant est persisté dans un index vectoriel aux côtés de
l'identifiant de l'activité. À la recherche, la requête de l'agent (typiquement
le sujet d'une activité à créer) est elle-même plongée, puis comparée à toutes
les entrées de l'index ; les $k$ plus proches sont retournées avec leur score
de similarité.

\subsection{Place de l'indexation dans l'architecture}

L'indexation n'est pas un service autonome que les agents appellent
directement. Elle est encapsulée dans le serveur \gls{mcp} sous la forme
d'outils de recherche sémantique typés, accompagnés d'une mise à jour
automatique de l'index à chaque création d'activité ou d'unité. L'agent ne
voit donc l'indexation qu'à travers deux interfaces : une opération de
recherche par similarité, en lecture, et une persistance enrichie de
l'indexation, en écriture.

Deux niveaux d'usage de cette indexation se distinguent dans la conception
du système.

Le premier niveau est l'usage \textbf{intra-syllabus}. Lorsqu'un agent doit
produire une nouvelle activité dans un syllabus existant, son premier réflexe
est la \textbf{lecture directe} des unités et activités sœurs : il interroge
le serveur \gls{mcp} pour lister la structure du syllabus, puis lit le corps
des activités voisines pour comprendre ce qui a déjà été enseigné. Cette
lecture de proximité, déterministe et peu coûteuse, porte l'essentiel de la
non-redondance dans le périmètre d'une seule unité. La recherche par
similarité y joue le rôle de \textbf{signal complémentaire} : elle permet de
détecter un quasi-doublon sémantique que la simple comparaison de titres
n'aurait pas révélé, ou de signaler qu'une activité jugée nouvelle reformule
en réalité un contenu existant.

Le second niveau est l'usage \textbf{global}, où la lecture directe ne peut
plus suffire. Trois situations le réclament. D'abord, la création d'un
\textbf{nouveau syllabus} : pour éviter qu'un enseignant ne produise un
doublon d'un syllabus existant dans la base, l'agent doit pouvoir interroger
l'ensemble du catalogue par similarité, et pas seulement la liste textuelle
des titres. Ensuite, l'\textbf{agent assistant élève} : lorsqu'un apprenant
est bloqué sur un concept, l'agent doit pouvoir retrouver les activités
sémantiquement proches à recommander, indépendamment de l'unité courante.
Enfin, l'\textbf{agent tuteur} : pour proposer des variantes pédagogiques ou
du matériel d'appui, l'agent doit pouvoir retrouver dans l'ensemble du
catalogue les activités voisines d'une activité donnée.

C'est sur ces trois usages que repose la valeur structurante de l'indexation.
Elle reste accessible au niveau intra-syllabus comme garde-fou anti-doublon,
mais son rôle décisif se déploie à l'échelle globale, dès que les agents
doivent raisonner au-delà du périmètre d'un seul syllabus.

%% ────────────────────────────────────────────────────────────
\section{Le protocole MCP appliqué au système}
\label{sec:mcp-systeme}

\subsection{Rappel : qu'est-ce que le MCP}

Le Model Context Protocol (\gls{mcp}), introduit en détail au chapitre~1, est
un protocole ouvert qui standardise la façon dont un agent \gls{llm} découvre
et appelle des outils fournis par un serveur externe~\cite{anthropic2024,mcpspec}.
Son architecture client--serveur repose sur \gls{json}-\gls{rpc} : le client
(intégré dans le runtime de l'agent) charge les descriptions d'outils, le
modèle décide lesquels appeler, et le serveur exécute la logique
correspondante puis renvoie un résultat structuré.

Deux propriétés du protocole sont particulièrement pertinentes pour
le présent système. La première est la \emph{découverte automatique} :
un agent conforme peut interroger un serveur \gls{mcp} pour obtenir
la liste de ses outils, y compris leurs schémas d'entrée et leur
documentation, ce qui permet d'ajouter un nouvel outil au serveur
sans modifier le code de l'agent — un simple redémarrage de la
connexion suffit. La seconde est le \emph{découplage entre l'agent
et son backend} : l'agent ne connaît ni la base de données ni le
langage du serveur, mais seulement des noms d'outils et des schémas
\gls{json}, ce qui permet de faire évoluer le backend (changer de
\gls{bd}, ajouter un cache) sans impacter les agents qui le
consomment.

\subsection{Architecture MCP du système}
\label{sec:mcp-archi}

Le système expose trois serveurs \gls{mcp}, chacun couvrant un domaine
fonctionnel distinct (figure~\ref{fig:mcp-servers}).

\begin{figure}[H]
\centering
\begin{tikzpicture}[
    server/.style={rectangle, rounded corners=5pt, draw, thick,
      minimum width=3.0cm, minimum height=1.0cm, align=center,
      fill=myblue!10, font=\small\bfseries},
    consumer/.style={rectangle, rounded corners=5pt, draw, thick,
      minimum width=3.0cm, minimum height=0.9cm, align=center,
      fill=gray!8, font=\small},
    arrow/.style={->, thick, -{Stealth[length=3mm]}, gray!70}
]
\node[server] (s1) {Serveur MCP\\Curriculum};
\node[server, right=0.8cm of s1] (s2) {Serveur MCP\\Données étudiants};
\node[server, right=0.8cm of s2] (s3) {Serveur MCP\\Analytique};

\node[consumer, below=1.5cm of s1] (c1)
  {Agents\\générateurs};
\node[consumer, below=1.5cm of s2] (c2)
  {Tuteur /\\Assistant élève};
\node[consumer, below=1.5cm of s3] (c3)
  {Assistant\\administrateur};

\draw[arrow] (s1) -- (c1);
\draw[arrow] (s2) -- (c2);
\draw[arrow] (s3) -- (c3);
\end{tikzpicture}
\caption{Trois serveurs MCP et les agents qui les consomment.}
\label{fig:mcp-servers}
\end{figure}

Le \emph{Serveur Curriculum} expose la hiérarchie syllabus
$\rightarrow$ unité $\rightarrow$ activité ; il constitue le serveur
principal du système, en cela que tous les agents générateurs le
consomment et que l'indexation vectorielle y est directement
intégrée. Le \emph{Serveur Données étudiants} expose les
soumissions, la progression et les retours ; il est consommé par le
tuteur, en lecture comme en écriture de feedback, et par
l'assistant élève, en lecture seule. Le \emph{Serveur Analytique}
expose enfin les rapports agrégés d'inscription, de complétion et
d'empreinte de contenu ; il est consommé par l'assistant
administrateur.

Chaque serveur supporte deux modes de transport conformes à la
spécification \gls{mcp} 2025~\cite{mcpspec}. En mode \emph{stdio},
le processus \gls{api} lance le serveur \gls{mcp} comme processus
enfant : ce mode est adapté au développement local, sans aucune
configuration réseau. En mode \emph{Streamable HTTP}, le serveur
\gls{mcp} tourne comme service indépendant derrière un domaine
privé ; chaque appel \gls{json}-\gls{rpc} prend la forme d'une
requête \gls{http} POST sans état, ce qui rend possible le passage
à l'échelle derrière un répartiteur de charge.

\subsection{Catalogue d'outils MCP}
\label{sec:mcp-catalogue}

Le tableau~\ref{tab:mcp-curriculum} énumère l'ensemble des outils exposés par
le serveur Curriculum. La distinction lecture / écriture est structurante :
les outils de lecture retournent des métadonnées légères (identifiants,
titres, ordre), tandis que les outils d'écriture effectuent exactement une
mutation et retournent la ligne insérée.

\begin{table}[H]
\centering
\caption{Outils MCP du serveur Curriculum.}
\label{tab:mcp-curriculum}
\small
\begin{tabularx}{\linewidth}{@{}>{\raggedright\arraybackslash}p{4.8cm}
    >{\raggedright\arraybackslash}p{1.6cm} X@{}}
\toprule
\textbf{Outil} & \textbf{Mode} & \textbf{Description}\\
\midrule
\mcptool{list_syllabuses}            & Lecture  & Liste des syllabus (id, titre, matière).\\
\mcptool{get_syllabus}               & Lecture  & Syllabus complet avec unités et activités.\\
\mcptool{list_unities}               & Lecture  & Unités d'un syllabus (id, titre, ordre).\\
\mcptool{list_activities_for_unity}  & Lecture  & Activités d'une unité (métadonnées légères).\\
\mcptool{get_activity}               & Lecture  & Corps complet d'une activité.\\
\mcptool{find_related_unities}       & Lecture  & Recherche sémantique d'unités voisines (par plongement).\\
\mcptool{find_related_activities}    & Lecture  & Recherche sémantique d'activités voisines (par plongement).\\
\addlinespace
\mcptool{create_syllabus}        & Écriture & Crée un syllabus.\\
\mcptool{create_unity}           & Écriture & Crée une unité dans un syllabus.\\
\mcptool{create_activity}        & Écriture & Crée une activité dans une unité (avec indexation automatique).\\
\mcptool{update_unity}           & Écriture & Met à jour le titre, l'ordre ou les objectifs d'une unité.\\
\mcptool{update_activity}        & Écriture & Met à jour le corps ou les métadonnées d'une activité.\\
\addlinespace
\multicolumn{3}{@{}l@{}}{\textit{Outils ajoutés pour la Version 2 (atelier formulaire) :}}\\
\mcptool{get_unit}                       & Lecture  & Lit une unité V2 (titre, brief, syllabus parent).\\
\mcptool{list_unit_activities}           & Lecture  & Liste les activités d'une unité filtrées par statut (clé de la règle anti-redite \og référencer, ne pas redire \fg{}).\\
\mcptool{create_unit_activity}           & Écriture & Persiste une activité V2 (\mcptool{markdown} + \mcptool{activity_json}).\\
\mcptool{update_unit_activity_status}    & Écriture & Fait transiter l'activité dans la machine d'état (figure~\ref{fig:v2-status}).\\
\bottomrule
\end{tabularx}
\end{table}

Les tableaux~\ref{tab:mcp-etudiants} et~\ref{tab:mcp-analytique} listent
respectivement les outils du serveur Données étudiants et du serveur
Analytique.

\begin{table}[H]
\centering
\caption{Outils MCP du serveur Données étudiants.}
\label{tab:mcp-etudiants}
\small
\begin{tabularx}{\linewidth}{@{}>{\raggedright\arraybackslash}p{4.8cm}
    >{\raggedright\arraybackslash}p{1.6cm} X@{}}
\toprule
\textbf{Outil} & \textbf{Mode} & \textbf{Description}\\
\midrule
\mcptool{get_submissions}       & Lecture  & Soumissions d'un étudiant (filtre optionnel par activité).\\
\mcptool{get_student_progress} & Lecture  & Progression d'un étudiant dans un syllabus.\\
\mcptool{save_feedback}         & Écriture & Persiste un retour correctif sur une soumission.\\
\bottomrule
\end{tabularx}
\end{table}

\begin{table}[H]
\centering
\caption{Outils MCP du serveur Analytique.}
\label{tab:mcp-analytique}
\small
\begin{tabularx}{\linewidth}{@{}>{\raggedright\arraybackslash}p{4.8cm}
    >{\raggedright\arraybackslash}p{1.6cm} X@{}}
\toprule
\textbf{Outil} & \textbf{Mode} & \textbf{Description}\\
\midrule
\mcptool{get_enrollment_report} & Lecture & Totaux inscriptions étudiants / syllabus.\\
\mcptool{get_completion_report} & Lecture & Nombre de soumissions et score moyen.\\
\mcptool{get_content_report}    & Lecture & Empreinte de contenu (syllabus, unités, activités).\\
\mcptool{get_user_stats}        & Lecture & Répartition par rôle.\\
\mcptool{list_students}          & Lecture & Liste des étudiants.\\
\bottomrule
\end{tabularx}
\end{table}

\subsection{Principe du moindre privilège}
\label{sec:moindre-privilege}

Chaque agent ne reçoit qu'un sous-ensemble des outils \gls{mcp}, strictement
nécessaire à son rôle. Cette discipline matérialise le \textbf{principe du
moindre privilège} : un agent qui n'a pas besoin d'écrire ne dispose d'aucun
outil d'écriture, ce qui élimine par construction toute mutation
accidentelle.

Le tableau~\ref{tab:whitelist} présente la liste blanche d'outils par agent.

\begin{table}[H]
\centering
\caption{Liste blanche d'outils MCP par agent.}
\label{tab:whitelist}
\small
\begin{tabularx}{\linewidth}{@{}>{\raggedright\arraybackslash}p{4.6cm} X@{}}
\toprule
\textbf{Agent} & \textbf{Outils autorisés}\\
\midrule
Générateur d'activités\newline(sans MCP) &
  Aucun outil \gls{mcp} : génère à partir de ses connaissances uniquement.\\
\addlinespace
Générateur d'activités\newline(avec MCP) &
  \mcptool{get_syllabus},\,
  \mcptool{list_unities},\,
  \mcptool{list_activities_for_unity},\,
  \mcptool{get_activity},\,
  \mcptool{find_related_activities},\,
  \mcptool{find_related_unities},\,
  \mcptool{create_activity},\,
  \mcptool{update_activity}\\
\addlinespace
Générateur de syllabus &
  \mcptool{list_syllabuses},\,
  \mcptool{get_syllabus},\,
  \mcptool{create_syllabus},\,
  \mcptool{list_unities},\,
  \mcptool{create_unity},\,
  \mcptool{update_unity}\\
\addlinespace
Tuteur &
  \mcptool{get_activity},\,
  \mcptool{list_activities_for_unity},\,
  \mcptool{find_related_activities},\,
  \mcptool{get_student_progress},\,
  \mcptool{get_submissions},\,
  \mcptool{save_feedback}\\
\addlinespace
Assistant élève &
  \mcptool{get_syllabus},\,
  \mcptool{list_unities},\,
  \mcptool{list_activities_for_unity},\,
  \mcptool{get_activity},\,
  \mcptool{find_related_activities},\,
  \mcptool{get_student_progress}\\
\addlinespace
Assistant administrateur &
  \mcptool{get_enrollment_report},\,
  \mcptool{get_completion_report},\,
  \mcptool{get_content_report},\,
  \mcptool{get_user_stats},\,
  \mcptool{list_students},\,
  \mcptool{list_syllabuses}\\
\bottomrule
\end{tabularx}
\end{table}

Trois observations méritent d'être soulignées. L'agent~1, qui sert
de générateur d'activités sans \gls{mcp}, ne dispose d'aucun outil
exposé par le protocole : sa production, non ancrée dans le contenu
existant, constitue la ligne de base à laquelle est comparée celle
de l'agent~2 dans la lecture comparative de la
\textsection\ref{sec:comparatif}. L'assistant élève, pour sa
part, ne dispose d'aucun outil d'écriture : il peut consulter le
syllabus et la progression de l'étudiant, mais ne peut ni créer de
contenu ni modifier de données. L'assistant administrateur, enfin,
n'a accès qu'à des outils de lecture agrégée : il ne voit jamais ni
le contenu individuel des activités, ni les soumissions détaillées,
ce qui matérialise concrètement le principe du moindre privilège au
niveau de la couche d'agents.

%% ────────────────────────────────────────────────────────────
\section{Catalogue des agents}
\label{sec:catalogue-agents}

La première itération de la couche d'agents livrée à la plateforme
SprintUp comportait cinq agents coordonnés par un orchestrateur unique
et une seule interface de création de syllabus, de nature
conversationnelle. Pour valider expérimentalement le compromis
\emph{ergonomie-temps de génération-contrôlabilité} à l'échelle de
SprintUp, nous avons développé trois variantes supplémentaires du même
agent~3 (création de syllabus). Ces quatre implémentations partagent le
même contrat fonctionnel mais reposent sur deux choix d'architecture
indépendants : (i)~le \emph{mode d'interaction utilisateur} et (ii)~la
\emph{stratégie d'orchestration} interne. Cette section présente d'abord
ces deux axes (\textsection\ref{sec:axes}), puis détaille les agents
communs aux deux modes d'interaction (générateurs d'activités, tuteur,
assistant élève, assistant administrateur), et termine sur l'agent~3
décliné en quatre variantes selon les six rubriques normalisées :
identité, outils, workflow, schéma de sortie, barre de qualité et
règles.

\subsection{Cadre d'analyse : modes d'interaction et stratégies
d'orchestration}
\label{sec:axes}

Les quatre variantes de l'agent~3 proposées à l'équipe SprintUp se
lisent comme le croisement de deux axes orthogonaux. Chaque axe répond
à une question d'intégration distincte ; nous justifions chacune
brièvement avant de présenter la matrice synthétique de la
figure~\ref{fig:matrice-agents}.

\paragraph{Axe~1 — Mode d'interaction utilisateur.}
Le premier axe désigne l'interface par laquelle un enseignant de
SprintUp déclenche la création d'un syllabus. Dans la Version~1, dite
\emph{interaction conversationnelle}, l'enseignant dialogue avec
l'agent dans un fil de discussion ; un protocole de streaming typé
lui livre les artefacts (plan, unités, leçons) au fur et à mesure
de leur production, et les questions de l'agent prennent la forme
d'\emph{interruptions} structurées (formulaire d'intake, demandes
de clarification). Cette variante adopte ainsi le paradigme
\emph{Human-in-the-Loop} (\gls{hitl}), qui autorise l'enseignant à
corriger une dérive en cours de génération. La Version~2, dite
\emph{interaction par atelier formulaire}, fait au contraire reposer
le déclenchement sur un formulaire structuré (titre, description,
audience, portée horaire, intention pédagogique) : l'enseignant crée
manuellement les unités et leurs activités sous forme de
\emph{brouillons}, puis déclenche explicitement la génération. Aucun
streaming de tokens ni interruption ne se produit en cours
d'exécution ; la progression d'une activité est uniquement exposée
par une machine d'état publique (figure~\ref{fig:v2-status}), ce qui
rend la surface consommable par un client externe sans avoir à
interpréter le protocole de streaming de la Version~1.

\paragraph{Axe~2 — Stratégie d'orchestration interne.}
Le second axe décrit la manière dont la couche d'agents organise,
en interne, les quatre phases canoniques d'un travail de
curation pédagogique — rechercher, planifier, rédiger et critiquer.
Dans l'\emph{orchestration monolithique} (Option~1), un graphe d'état
figé enchaîne ces phases dans un ordre prédéfini et un unique
superviseur prend toutes les décisions de routage : le flot y gagne
en lisibilité et en reproductibilité, au prix d'une moindre
spécialisation par phase. Dans l'\emph{orchestration par sous-agents
spécialistes} (Option~2), un superviseur délègue au contraire chaque
phase à un sous-agent porteur d'une instruction système dédiée
(planner, writer, activity-maker, critic) ; la communication entre
sous-agents passe alors soit par un \gls{vfs} partagé (canal
documentaire), soit par un ordonnanceur \textsc{TypeScript} pur
(canal d'invocation). Le flot gagne en spécialisation par rôle, mais
au prix d'une coordination plus riche à mettre en place et à
maintenir.

\paragraph{Synthèse : matrice des quatre variantes.}
La figure~\ref{fig:matrice-agents} présente le croisement des deux
axes ; les quatre variantes apparaissent comme les cellules du
tableau, identifiées par le nom de code de leur composant principal
dans le dépôt : \mcptool{syllabus-generator} (V1/Op1),
\mcptool{deepagent} (V1/Op2), \mcptool{unit-bulk-supervisor} associé à
\mcptool{unit-activity-generator} (V2/Op1) et
\mcptool{sdk-orchestrator} (V2/Op2). Les quatre sous-sections de
l'agent~3 (\textsection\ref{sec:agent3}) détaillent ensuite chacune
de ces variantes selon les six rubriques normalisées.

\begin{figure}[H]
\centering
\begin{tikzpicture}[
  font=\small,
  colhead/.style={font=\small\bfseries, align=center,
                  text width=5.4cm},
  rowhead/.style={font=\small\bfseries, align=center,
                  text width=2.8cm, rotate=90},
  cell/.style ={rectangle, rounded corners=4pt, draw=myblue!60,
                line width=0.6pt, fill=myblue!5,
                minimum width=5.6cm, minimum height=2.9cm,
                text width=5.2cm, inner sep=4pt,
                align=center, font=\footnotesize},
]

%% --- axe horizontal (modes d'interaction) ---
\node[colhead] (h1) at (0, 3.25) {V1 — Interaction\\conversationnelle};
\node[colhead] (h2) at (6.4, 3.25) {V2 — Interaction par\\atelier formulaire};

%% --- axe vertical (stratégie d'orchestration) ---
\node[rowhead] (r1) at (-4.6, 0.6)  {Op.~1\\Orchestration\\monolithique};
\node[rowhead] (r2) at (-4.6, -2.7) {Op.~2\\Sous-agents\\spécialistes};

%% --- cellules ---
\node[cell] (c11) at (0, 0.6)
  {{\bfseries\color{myblue!80!black}\texttt{syllabus-generator}}\\[2pt]
   superviseur à cinq actions\\
   (intake, search, ask, write, reply)\\[2pt]
   sous-graphe \emph{search} parallèle\\
   + portail critique sévérité-sensible};

\node[cell] (c12) at (6.4, 0.6)
  {{\bfseries\color{myblue!80!black}\texttt{sdk-orchestrator}}\\[2pt]
   planificateur \textsc{LLM}\\
   $\rightarrow$ fan-out \texttt{Send}\\
   $\rightarrow$ rédacteur par unité\\[2pt]
   unités en parallèle,\\
   activités en séquence};

\node[cell] (c21) at (0, -2.7)
  {{\bfseries\color{myblue!80!black}\texttt{deepagent}}\\[2pt]
   superviseur \emph{deepagents}\\
   + 4 sous-agents : planner /\\
   writer / activity-maker / critic\\[2pt]
   communication par\\
   \textsc{VFS} partagé};

\node[cell] (c22) at (6.4, -2.7)
  {{\bfseries\color{myblue!80!black}\texttt{unit-bulk-supervisor}}\\[2pt]
   + \texttt{unit-activity-generator}\\
   + \texttt{unit-suggestion},\\
   \texttt{activity-suggestion}\\[2pt]
   brouillons +\\
   vagues de dépendances};
\end{tikzpicture}
\caption{Quatre variantes de l'agent~3 proposées pour SprintUp,
  positionnées au croisement du mode d'interaction utilisateur
  (Version~1 vs.\ Version~2) et de la stratégie d'orchestration
  (monolithique vs.\ sous-agents spécialistes). Chaque cellule porte
  le nom de code de la variante dans le dépôt et résume son schéma
  d'orchestration.}
\label{fig:matrice-agents}
\end{figure}

\paragraph{Canaux d'inter-communication.}
Les trois schémas d'orchestration distincts font apparaître trois
canaux d'inter-communication entre composants, dont la nature
conditionne directement le coût d'intégration dans SprintUp.
En V1/Op.~1 (\mcptool{syllabus-generator}), les agents ne
communiquent pas directement entre eux : toute coordination passe
par l'orchestrateur \textsc{LLM} unique, qui décide à chaque tour
quelle action exécuter et porte donc, à lui seul, l'intégralité de
la logique de routage. En V1/Op.~2 (\mcptool{deepagent}), c'est un
\gls{vfs} partagé qui sert de canal documentaire entre sous-agents
(\mcptool{/pedagogy\_plan.md}, \mcptool{/lessons/<id>.md},
\mcptool{/activities/<id>.json}, \mcptool{/critiques/<target>.md})~:
chaque sous-agent lit le travail des précédents et y ajoute le sien,
ce qui rend les artefacts inspectables a posteriori. En Version~2
enfin, il n'existe plus d'orchestrateur \textsc{LLM} de niveau
\emph{supervisor} : l'ordonnancement est intégralement délégué à du
code \textsc{TypeScript} — soit le \mcptool{unit-bulk-supervisor}
de la V2/Op.~1, qui exécute les activités par vagues de dépendances,
soit l'orchestrateur \gls{sdk} de la V2/Op.~2, qui exécute les
unités en parallèle au moyen de la primitive \mcptool{Send} de
LangGraph. Cette distinction n'est pas qu'esthétique : elle dicte si
SprintUp peut auditer l'orchestration via du code (Version~2) ou
seulement via des traces \textsc{LLM} (Version~1).

\subsection{Agent 1 : Générateur d'activités sans MCP}
\label{sec:agent1}

\paragraph{Identité.} Agent de génération d'activités pédagogiques pour la
plateforme SprintUp. Produit une ou plusieurs activités (cours, question,
exercice ou quiz) pour un sujet et une audience donnés, \emph{sans consulter
le contenu existant}. Cet agent constitue la ligne de base du système : son
fonctionnement repose entièrement sur les connaissances du modèle de langue
sous-jacent et sur les paramètres fournis dans la requête.

\paragraph{Outils.} Aucun outil \gls{mcp}, aucune recherche documentaire
externe, aucune lecture de la base de données de la plateforme. L'agent
répond à partir du seul prompt utilisateur. Cette absence délibérée en fait
le référent de comparaison pour apprécier la valeur ajoutée du \gls{mcp} et
de l'indexation au niveau de l'agent~2.

\paragraph{Workflow.}
Le déroulé est simple en quatre temps. L'agent commence par
identifier dans la requête l'ensemble des paramètres utiles — sujet,
niveau d'audience, type d'activité souhaité, nombre d'éléments,
difficulté cible et langue. Il planifie ensuite mentalement la
courbe de difficulté sur l'ensemble des activités demandées, de
manière à éviter une croissance monotone ou un palier inutile.
Vient alors la rédaction proprement dite : chaque activité est
rédigée en variant délibérément sa formulation et son niveau de
Bloom par rapport à la précédente, afin d'éviter la répétition.
L'agent termine en retournant un tableau \gls{json} dont chaque
objet décrit une activité complète.

\paragraph{Schéma de sortie.} Tableau \gls{json} d'objets \texttt{Activity} :
\mcptool{title}, \mcptool{activity_type} (\texttt{course}, \texttt{question},
\texttt{exercise}, \texttt{quiz}), \mcptool{content}, \mcptool{description},
\mcptool{difficulty} (1--5), \mcptool{duration_min}, \mcptool{bloom_level},
\mcptool{learning_objectives}, \mcptool{prerequisites}, \mcptool{key_terms},
\mcptool{mcqs} (pour les quiz : question, 4~options, \mcptool{correct_index},
explication).

\paragraph{Barre de qualité.} Une activité produite par cet agent
doit présenter un vocabulaire, des exemples et une complexité alignés
sur le niveau de l'audience cible. Chaque activité doit comporter au
moins un objectif d'apprentissage explicite, et les exemples résolus
doivent inclure les étapes intermédiaires plutôt que de se limiter
à la réponse finale. Pour les quiz, les distracteurs doivent être
plausibles, faute de quoi la difficulté réelle se trouve supérieure
à ce que le \mcptool{difficulty} déclare.

\paragraph{Règles.} Quatre interdits encadrent l'agent. Il ne doit
pas générer plus ni moins d'activités que demandé dans la requête.
Il ne doit pas répondre avec autre chose que le tableau \gls{json}
attendu — ni texte d'accompagnement, ni explication libre. Il ne
doit pas inclure de données personnelles, de noms réels de mineurs
ni de texte sous droit d'auteur. Enfin, il ne doit faire aucune
hypothèse sur le syllabus ou les activités existantes : par
construction, l'agent travaille en aveugle, et c'est précisément ce
qui en fait la ligne de base de comparaison pour l'agent~2.

\subsection{Agent 2 : Générateur d'activités avec MCP et indexation}
\label{sec:agent2}

\paragraph{Identité.} Version enrichie de l'agent~1. Avant de produire du
contenu, il consulte le syllabus et l'unité existants via \gls{mcp} afin que
les nouvelles activités (a)~s'inscrivent dans la progression pédagogique et
(b)~ne dupliquent pas le matériel existant. La non-redondance s'appuie
d'abord sur la lecture directe des activités sœurs de l'unité cible et,
en complément, sur la recherche par similarité mise à disposition par
l'indexation.

\paragraph{Outils.} L'agent dispose, en lecture sur le serveur MCP
Curriculum, des outils nécessaires pour situer la requête dans le
syllabus existant et y repérer le risque de redondance. Il peut
lire la hiérarchie complète d'un syllabus
(\mcptool{get_syllabus(syllabus_id)}), énumérer ses unités
(\mcptool{list_unities(syllabus_id)}) et les activités déjà
présentes dans l'unité cible
(\mcptool{list_activities_for_unity(unity_id)}), puis ouvrir le
corps d'une activité particulière (\mcptool{get_activity(activity_id)}).
Il mobilise enfin l'indexation vectorielle au travers de deux outils
de recherche sémantique :
\mcptool{find_related_activities(syllabus_id, query)} pour repérer
les activités voisines d'un sujet, et
\mcptool{find_related_unities(syllabus_id, query)} pour situer la
demande dans la structure des unités. En écriture, l'agent dispose
de \mcptool{create_activity(...)}, qui persiste une nouvelle
activité et la réindexe automatiquement, et de
\mcptool{update_activity(...)}, qu'il privilégie dès que la lecture
des voisines révèle une activité à enrichir plutôt qu'une nouvelle
à créer.

\paragraph{Workflow.}
L'agent commence par extraire de la requête les paramètres
structurants (sujet, \mcptool{syllabus_id}, \mcptool{unity_id},
nombre d'activités, difficulté cible). Il procède alors à une
\emph{phase de lecture} dont l'objectif est double : confirmer
l'existence et inspecter la structure globale du syllabus via
\mcptool{get_syllabus(syllabus_id)}, puis lire les activités déjà
présentes dans l'unité cible via
\mcptool{list_activities_for_unity(unity_id)}, en ouvrant le corps
de celles qui paraissent les plus proches à l'aide de
\mcptool{get_activity}. Cette lecture de proximité est, en pratique,
le signal principal d'anti-redondance ; elle est complétée par un
appel sémantique à \mcptool{find_related_activities(syllabus_id, query)}
qui vérifie qu'aucune activité proche n'existe ailleurs dans le
syllabus.

Le résultat de cette phase de lecture oriente alors la décision. Si
une activité sœur couvre déjà le sujet, l'agent répond par un objet
\gls{json} décrivant le conflit ou propose une mise à jour ciblée
via \mcptool{update_activity}, plutôt que de créer un doublon
déguisé. Sinon, il planifie la courbe de difficulté sur les
activités demandées, puis rédige chacune d'elles en l'inscrivant
explicitement dans la progression de l'unité. Pour chaque activité
rédigée, un appel à \mcptool{create_activity} en assure la
persistance et déclenche, de façon transparente, la mise à jour de
l'index vectoriel. L'agent retourne enfin un tableau \gls{json}
comportant l'ensemble des enregistrements persistés avec leur
identifiant.

\paragraph{Schéma de sortie.} En cas de succès : tableau \gls{json} des
enregistrements retournés par \mcptool{create_activity} (avec \texttt{id}).
En cas de doublon détecté :
\begin{center}
\texttt{\{"skipped": true,\,"reason": "duplicate detected",\,"existing": [...]\}}
\end{center}

\paragraph{Barre de qualité.} Trois exigences guident l'évaluation
d'une activité produite par cet agent. Elle doit d'abord apporter
quelque chose que l'unité n'avait pas encore, faute de quoi elle
n'a aucune raison d'exister. Sa difficulté doit ensuite s'inscrire
dans la progression de l'unité — ni rupture brutale, ni régression
inexpliquée. Sa cohérence pédagogique, enfin, suppose que les
prérequis qu'elle invoque renvoient bien à des concepts déjà
enseignés plus tôt dans le syllabus, ce que la lecture en amont
permet de vérifier.

\paragraph{Règles.} L'agent obéit à trois interdits stricts. Il ne
doit \emph{jamais} appeler \mcptool{create_activity} avant d'avoir
lu les activités sœurs de l'unité cible : l'ordre lecture-puis-écriture
est la condition nécessaire de l'anti-redondance. Il ne doit
\emph{jamais} créer une activité quand un quasi-doublon est
identifié, que la détection provienne de la lecture directe des
voisines ou de la recherche sémantique. Il ne doit enfin
\emph{jamais} fabriquer un \mcptool{unity_id} ou un
\mcptool{syllabus_id} qui n'aurait pas été fourni dans la requête :
en leur absence, l'agent demande une clarification.

\subsection{Agent 3 : Création de syllabus — quatre variantes
implémentées}
\label{sec:agent3}

La création d'un syllabus est le seul agent dont nous avons développé
plusieurs implémentations à destination de SprintUp, afin de couvrir
les deux axes présentés en \textsection\ref{sec:axes}. Cette
sous-section décrit les quatre variantes effectivement livrées, dans
l'ordre suivant : V1/Op.~1 (\texttt{syllabus-generator}, conversation +
\gls{hitl}), V1/Op.~2 (\texttt{deepagent}, conversation + sous-agents
spécialistes), V2/Op.~1 (atelier manuel,
\texttt{unit-bulk-supervisor}) et V2/Op.~2
(\texttt{sdk-orchestrator}). L'identité commune aux quatre variantes
est celle d'un \emph{concepteur de curriculum senior} qui produit la
structure d'un parcours pédagogique complet (syllabus + unités +
activités) à partir d'un sujet, d'un profil d'audience et d'une
enveloppe horaire ou d'un objectif terminal. Les variantes se
distinguent par leur mode d'interaction (conversation vs.\ atelier),
par leur granularité de production (structure seule vs.\ structure +
contenu d'activité) et par leur stratégie d'orchestration interne
(monolithique vs.\ sous-agents spécialistes).

%% ---------------------------------------------------------
\subsubsection{V1 / Option~1 — \texttt{syllabus-generator} (HITL)}
\label{sec:v1-op1}

\paragraph{Identité.} Agent conversationnel construit autour d'un
superviseur LangGraph qui classe chaque tour utilisateur parmi cinq
actions — \mcptool{intake}, \mcptool{search}, \mcptool{ask},
\mcptool{write}, \mcptool{reply} — et d'un sous-graphe
\emph{rechercher / écrire / critiquer} qui matérialise une bonne partie
de la valeur pédagogique. La conversation tient lieu de mémoire de
thread ; les indécisions du superviseur sont résolues par des
\emph{interruptions} structurées (\emph{Human-in-the-Loop}).

\paragraph{Outils.} \textbf{Lecture (MCP Curriculum)} :
\mcptool{list_syllabuses}, \mcptool{get_syllabus},
\mcptool{list_unities}, \mcptool{find_related_unities},
\mcptool{find_related_activities}. \textbf{Écriture (MCP Curriculum)} :
\mcptool{create_syllabus}, \mcptool{create_unity},
\mcptool{update_unity}, \mcptool{create_activity}. \textbf{Recherche
externe} : Serper (web), assortie d'un scraper et d'un summarizer dans
le sous-graphe \mcptool{search}.

\paragraph{Workflow.}
À chaque tour de conversation, le superviseur commence par classer
le message utilisateur parmi les cinq actions évoquées plus haut.
Lorsque la requête initiale manque d'informations structurantes —
niveau d'audience, durée totale, objectif terminal, prérequis ou
langue — le superviseur interrompt le graphe et émet un formulaire
typé (\mcptool{intake_form}). La reprise injecte les valeurs renseignées
comme contraintes dures dans l'état (\mcptool{audience.level},
\mcptool{scope.duration_hours}, \mcptool{scope.target_outcome},
\mcptool{audience.language}), de sorte que les phases ultérieures n'ont
plus à deviner ces paramètres.

Vient ensuite la phase de recherche, matérialisée par le sous-graphe
\mcptool{search} qui enchaîne quatre étapes (Serper par sujet, picker,
scraper, summarizer) parallélisées à l'aide de la primitive
\mcptool{Send} de LangGraph : une branche est lancée par sujet de
recherche, ce qui permet de couvrir plusieurs angles d'un même
curriculum sans sérialiser les appels web. Au sortir de cette phase,
le superviseur peut émettre au plus deux questions de clarification
post-recherche (\mcptool{ask}) si, et seulement si, le résumé fait
apparaître une ambiguïté concrète qui change la forme du syllabus.

La phase de planification consolide ces informations en une
structure complète : appel préalable à \mcptool{list_syllabuses}
pour garantir l'idempotence, choix de quatre à huit unités ordonnées
par chaîne de prérequis, et attribution d'objectifs pédagogiques
(taxonomie de Bloom) à chaque unité. La rédaction s'effectue ensuite
en \emph{vagues topologiques} : après un premier \mcptool{create_syllabus},
chaque vague rassemble toutes les leçons dont les dépendances ont
déjà été \emph{committées}, et leur écriture s'exécute en parallèle.
Chaque vague étend ainsi une clôture transitive de la relation
\mcptool{depends_on}. Pour chaque leçon écrite, un portail critique
\emph{sévérité-sensible} fait tourner la critique au plus une fois~:
en cas d'échec, le \emph{writer} révise une fois la leçon, qui est
alors committée en l'état. L'interaction se clôt par le retour, à
l'utilisateur, de la structure complète en \gls{json} (schéma
ci-dessous).

\paragraph{Schéma de sortie.} L'agent retourne un objet \gls{json}
comprenant l'identifiant et les métadonnées du syllabus
(\mcptool{syllabus_id}, \mcptool{title}, \mcptool{subject},
\mcptool{description}, \mcptool{duration_hours}), un tableau
\mcptool{units[]} dont chaque élément porte un \mcptool{id}, un
\mcptool{title}, un \mcptool{order} et la liste de ses
\mcptool{learning_objectives}, et, le cas échéant, un tableau
\mcptool{lessons[]} — le contenu produit en vagues figurant
ensuite dans le tableau \mcptool{activities[]} via les outils
d'écriture du serveur MCP.

\paragraph{Barre de qualité.} Un syllabus accepté doit présenter
une progression logique, au sens où les prérequis d'une unité sont
effectivement couverts par les unités précédentes. Chaque unité
doit en outre porter au moins un objectif Bloom du niveau
«~appliquer~» ou supérieur, sous peine de tomber dans la simple
récitation. Le cumul des durées des unités doit être cohérent avec
l'enveloppe horaire fixée par l'\mcptool{intake_form}, et la langue
des intitulés doit s'aligner sur la langue de l'audience cible.

\paragraph{Règles propres à cette voie.} Cette voie obéit à quatre
règles d'usage strictes. D'abord, lister avant de créer : tout appel
d'écriture est systématiquement précédé d'une lecture de l'existant
pour garantir l'idempotence. Ensuite, l'intake HITL est plafonné à
un unique \mcptool{intake_form} par fil de conversation, et les
questions de clarification (\mcptool{ask}) à deux occurrences
post-recherche, afin de préserver l'expérience conversationnelle.
La critique reste en lecture seule : la révision est confiée au
\emph{writer}, jamais au critic, ce qui clarifie les responsabilités
de chaque agent. Enfin, l'outil \mcptool{find_related_unities} est
mobilisé deux fois — à l'échelle du syllabus en cours puis à
l'échelle globale — afin de prévenir aussi bien les doublons intra-
qu'inter-syllabus.

%% ---------------------------------------------------------
\subsubsection{V1 / Option~2 — \texttt{deepagent}}
\label{sec:v1-op2}

\paragraph{Identité.} Superviseur \emph{deepagents} généraliste qui
décide au runtime quelle capacité utiliser pour chaque demande
(construire un syllabus, produire un worksheet, critiquer un artefact,
ou simplement répondre). La construction de syllabus résulte d'un
enchaînement de délégations \mcptool{task(name, briefing)} vers quatre
sous-agents spécialistes : \mcptool{pedagogy_planner},
\mcptool{writer}, \mcptool{activity_maker}, \mcptool{pedagogy_critic}.

\paragraph{Outils.} La trousse outillage du superviseur tient en
trois familles. La première est l'outil de délégation
\mcptool{task(name, briefing)}, qui invoque un sous-agent dans un
contexte \textsc{LLM} propre, c'est-à-dire sans héritage de
l'historique du superviseur. La deuxième est la liste blanche d'outils
MCP, distincte pour chaque sous-agent : le \mcptool{pedagogy_planner}
se contente de lectures (\mcptool{list_syllabuses}, \mcptool{get_syllabus}),
le \mcptool{writer} dispose des outils d'écriture du serveur Curriculum
(\mcptool{create_unity}, \mcptool{create_activity},
\mcptool{update_activity}) et l'\mcptool{activity_maker} se concentre sur
la production du worksheet structuré qu'il persiste également via
\mcptool{create_activity} ; le \mcptool{pedagogy_critic} reste
en lecture seule. Vient enfin le canal d'inter-communication,
matérialisé par un \gls{vfs} partagé : les sous-agents y échangent
leurs artefacts via les chemins \mcptool{/user_profile.md},
\mcptool{/pedagogy_plan.md}, \mcptool{/lessons/<id>.md},
\mcptool{/activities/<id>.json} et \mcptool{/critiques/<target>.md}.

\paragraph{Workflow.}
Lorsqu'une demande de construction de syllabus arrive, le superviseur
la classe comme telle et délègue d'abord au \mcptool{pedagogy_planner},
qui produit le plan pédagogique \mcptool{/pedagogy_plan.md} dans le
\gls{vfs}. Pour chaque chapitre de ce plan, le superviseur lance
ensuite un \mcptool{task} vers le \mcptool{writer} : celui-ci rédige
la leçon correspondante à \mcptool{/lessons/<id>.md} et la persiste
via les outils d'écriture MCP appropriés. De façon optionnelle,
l'\mcptool{activity_maker} produit un worksheet structuré (cours +
questions à choix multiples) et le \mcptool{pedagogy_critic} émet,
en lecture seule, des \emph{findings} tagués par niveau de sévérité.
Le superviseur synthétise enfin l'ensemble et retourne une carte UI
typée \mcptool{<artifact kind="syllabus" />} à l'utilisateur.

\paragraph{Schéma de sortie.} La sortie principale est une carte UI
typée \mcptool{<artifact kind="syllabus" />} décrivant la structure
créée ; le cas échéant s'y ajoutent une carte
\mcptool{<artifact kind="worksheet" />} et les \emph{findings}
émis par le \mcptool{pedagogy_critic}.

\paragraph{Barre de qualité.} La barre de qualité reprend les
critères de progression et de Bloom de la V1/Op.~1 (logique de
prérequis, au moins un objectif d'application par unité, durées
cohérentes avec l'enveloppe horaire) et y ajoute deux contraintes
spécifiques à la délégation. D'une part, le critic en lecture seule
ne modifie jamais le contenu : il se borne à le signaler, ce qui
limite les boucles de réécriture en cascade. D'autre part, chaque
appel \mcptool{task} part avec un contexte LLM vide : le briefing
fourni par le superviseur doit donc être suffisamment riche pour
que le sous-agent exécute sa tâche en autonomie (principe de
\emph{subagent context isolation}).

\paragraph{Règles propres à cette voie.} Trois règles d'usage
encadrent l'orchestration par sous-agents. L'isolation est stricte :
aucun contexte LLM n'est hérité d'un \mcptool{task} à l'autre, ce
qui garantit la reproductibilité et limite les fuites d'historique
entre spécialités. La séparation des rôles est tout aussi nette :
le planner n'écrit pas de leçon, le writer ne crée pas de syllabus,
l'activity-maker n'émet pas de plan global. Enfin, l'inter-communication
passe exclusivement par le \gls{vfs} : aucune mémoire RAM
partagée n'est exploitée entre sous-agents, ce qui préserve la
traçabilité de bout en bout et autorise un audit a posteriori des
artefacts échangés.

%% ---------------------------------------------------------
\subsubsection{V2 / Option~1 — atelier manuel
  (\texttt{unit-bulk-supervisor})}
\label{sec:v2-op1}

\paragraph{Identité.} Trois agents complémentaires opèrent sur la
surface atelier : \mcptool{unit-suggestion-agent} et
\mcptool{activity-suggestion-agent} (\emph{optionnels}, propositions en
sortie structurée à valider par l'utilisateur) et
\mcptool{unit-activity-generator} (génération du contenu d'une
activité). Un ordonnanceur \textsc{TypeScript} pur,
\mcptool{unit-bulk-supervisor}, coordonne les exécutions par vagues de
dépendances.

\paragraph{Outils MCP nouveaux pour V2.}
\mcptool{get_unit(unit_id)},
\mcptool{list_unit_activities(unit_id, status_in, ranking,
  token_budget)} (\emph{outil clé} pour la règle anti-redite),
\mcptool{create_unit_activity(...)},
\mcptool{update_unit_activity_status(...)}. Outils réutilisés de la
V1 : \mcptool{get_syllabus}, \mcptool{create_syllabus}.

\paragraph{Workflow.}
L'enseignant commence par créer, via le formulaire d'atelier, le
syllabus puis ses unités, et y dépose chaque activité sous forme de
\emph{brouillon} — un titre et une courte description qui servent
de germe au \mcptool{target_outcome} de la future activité. Au
besoin, il s'appuie sur les deux agents de suggestion :
\mcptool{unit-suggestion-agent} propose trois à quinze unités
cohérentes à partir du brief du syllabus, et
\mcptool{activity-suggestion-agent} propose deux à douze activités
par unité. Ces propositions sont renvoyées en sortie structurée et restent à
l'état de candidates : l'enseignant les accepte, les édite ou les
rejette avant qu'elles ne soient persistées comme brouillons.

La génération d'une activité individuelle suit ensuite la machine
d'état de la figure~\ref{fig:v2-status}. L'agent
\mcptool{unit-activity-generator} lit d'abord les activités voisines
déjà \mcptool{ready} via \mcptool{list_unit_activities}, en les
classant par récence et sous un budget de tokens borné, ce qui
implémente concrètement la règle anti-redite (\emph{no-repeat}) au
sein de l'unité. Il appelle ensuite le \textsc{LLM} \emph{writer} en
\emph{structured output} contre le schéma \mcptool{UnitActivityContract},
puis persiste les deux moitiés de l'activité (\mcptool{markdown} pour
la leçon, \mcptool{activity_json} pour le devoir) avant de faire
transiter le statut vers \mcptool{ready}. Lorsque l'enseignant
déclenche une génération en lot,
\mcptool{unit-bulk-supervisor} calcule les vagues de dépendances
définies par la relation \mcptool{depends_on} et lance jusqu'à $n$
activités en parallèle par vague (par défaut $n = 4$, avec un plafond
de seize), de manière à maximiser le débit sans rompre l'ordre
pédagogique.

\paragraph{Schéma de sortie (activité unique).} Chaque activité
générée possède un \mcptool{title} et un \mcptool{target_outcome}
exprimé par un verbe d'action de la taxonomie de Bloom. Son contenu
est scindé en deux moitiés. La \emph{moitié leçon}, portée par le
champ \mcptool{markdown}, prend la forme d'une prose structurée en
sous-titres H2 et H3 avec exemples et une section «~Try it~». La
\emph{moitié devoir}, portée par le champ \mcptool{activity_json},
est un objet typé \mcptool{type: "mcq_set"} accompagné d'un
\mcptool{intro?} optionnel et d'un tableau \mcptool{questions[]}~;
chaque question comprend un \mcptool{prompt}, quatre \mcptool{options[4]},
un \mcptool{correct_index} et une \mcptool{explanation?} facultative.
C'est ce même format qu'attend la plateforme SprintUp dans sa version
courante ; il reste extensible à d'autres types d'évaluation
(appariement, ordonnancement) en élargissant le champ \mcptool{type}.

\paragraph{Barre de qualité.} Deux exigences encadrent la qualité
d'une activité V2. La première est l'auto-suffisance : la moitié
devoir doit être \emph{répondable} à partir de la seule moitié
leçon, sans recours à une activité extérieure. La seconde est la
règle dite «~référencer, ne pas redire~» : il est encouragé de
citer les activités voisines pour ancrer la nouvelle dans son
contexte, mais il est strictement interdit de reprendre un énoncé,
un texte d'option ou un exemple déjà résolu. Cette règle est ce
qui transforme l'unité en \emph{parcours} et non en simple
collection.

\paragraph{Règles propres à cette voie.} Cette voie ne pratique
ni interruption \emph{Human-in-the-Loop} ni streaming de tokens :
toute la progression observable côté client passe par la colonne
\mcptool{unit_activities.status}, dont les transitions sont
décrites par la figure~\ref{fig:v2-status}. Toutes les écritures
en base de données franchissent par ailleurs la liste blanche
d'outils du \mcptool{unit-activity-generator}, ce qui borne
strictement les effets de bord et permet à l'équipe SprintUp
d'auditer chaque mutation a posteriori.

%% ---------------------------------------------------------
\subsubsection{V2 / Option~2 — \texttt{sdk-orchestrator}}
\label{sec:v2-op2}

\paragraph{Identité.} Graphe LangGraph à deux nœuds
(\mcptool{planner}, \mcptool{writer}) exposé en endpoint REST unique
(\mcptool{POST /api/v2/syllabuses/sdk-generate}). Réutilise l'agent
d'activité de la voie V2/Op1. Conçu pour être consommé par un
\gls{sdk} externe sans dialogue conversationnel.

\paragraph{Outils.} La voie V2/Op.~2 réutilise exactement la même
liste blanche d'outils MCP que la V2/Op.~1 (\mcptool{get_syllabus},
\mcptool{get_unit}, \mcptool{list_unit_activities},
\mcptool{create_unit_activity}, \mcptool{update_unit_activity_status}),
ce qui garantit la cohérence des règles d'écriture entre les deux
voies de la Version~2. Elle y ajoute un seul outil propre :
la primitive de fan-out \mcptool{Send} fournie par LangGraph, qui
permet à un nœud d'état d'émettre $n$ exécutions parallèles d'un
autre nœud, chacune avec son propre fragment d'état.

\paragraph{Workflow.}
Le flot est composé de deux nœuds enchaînés. Le premier, le
\emph{planner}, commence par lire le syllabus via MCP
(\mcptool{title}, \mcptool{audience}, \mcptool{scope},
\mcptool{pedagogy}), puis appelle le \textsc{LLM} superviseur en
\emph{structured output} contre le contrat
\mcptool{CoursePlanContract} (trois à quinze unités, deux à six
activités par unité). L'ensemble des brouillons d'unité et
d'activité produits est ensuite inséré en une seule transaction
(\emph{batch-insert}), de sorte qu'à ce stade la structure complète
existe déjà en base, même si aucun contenu n'a été écrit. Le
deuxième nœud, le \emph{writer}, est invoqué en fan-out par
\mcptool{state.unit_ids.map(uid => new Send("writer", \{unit_id\}))} :
une exécution writer est créée par unité, et ces exécutions tournent
en parallèle. À l'intérieur de chaque exécution writer, les
activités sont itérées dans l'ordre du champ \mcptool{order_index}
et générées séquentiellement par
\mcptool{unit-activity-generator.generate} — le même composant
qu'en V2/Op.~1, ce qui assure que la règle «~référencer, ne pas
redire~» est appliquée de façon identique.

L'effet net est celui décrit dans la cellule V2/Op.~2 de la
figure~\ref{fig:matrice-agents} : les unités sont parallèles entre
elles, mais les activités d'une même unité restent séquentielles.
Cette décomposition n'est pas anodine : la séquentialité intra-unité
préserve le parcours pédagogique linéaire — chaque activité voit
bel et bien, via MCP, les \mcptool{ready} qui la précèdent dans son
unité — tandis que la parallélisation inter-unités comprime
fortement le temps total de génération d'un syllabus complet.

\paragraph{Schéma de sortie.} Au niveau de chaque activité, le
schéma de sortie est rigoureusement identique à celui de la V2/Op.~1
(moitié leçon en \mcptool{markdown} et moitié devoir en
\mcptool{activity_json}). La progression globale du syllabus est
quant à elle exposée côté client par \emph{polling} de l'endpoint
\mcptool{GET /api/v2/syllabuses/:id}, à une cadence de deux secondes
côté front-end SprintUp.

\paragraph{Barre de qualité.} La voie V2/Op.~2 ajoute deux exigences
de qualité à celles déjà décrites pour les activités en V2/Op.~1.
D'abord, le contrat de cours produit par le \emph{planner} doit être
\emph{complet} avant la première écriture de contenu : toutes les
unités et toutes les activités sont définies en même temps, et
aucune activité n'est jamais rédigée avec une vue partielle de la
structure. Ensuite, trois invariants doivent être simultanément
maintenus : le parallélisme inter-unités, qui apporte le gain de
temps ; la séquentialité intra-unité, qui préserve la linéarité du
parcours~; et la règle «~référencer, ne pas redire~» héritée de la
V2/Op.~1, qui prévient la redondance à l'échelle de l'unité.

\paragraph{Règles propres à cette voie.} L'endpoint est délibérément
non bloquant : il répond \textsc{HTTP}~202 et la génération s'effectue
en arrière-plan, ce qui rend l'intégration dans un \gls{sdk} externe
immédiatement consommable. Le client n'observe l'avancement
qu'à travers la machine d'état publique (figure~\ref{fig:v2-status}),
sans aucune autre source de vérité. Enfin, le \emph{planner} n'a pas
le droit de mettre à jour une unité déjà écrite : toute modification
post-génération doit repasser par la surface atelier de la V2/Op.~1,
de manière à éviter qu'un appel SDK ne réécrive silencieusement du
contenu déjà validé par l'enseignant.

%% ---------------------------------------------------------
\subsubsection{Machine d'état d'une activité V2}
\label{sec:v2-status-text}

Les deux voies de la Version~2 partagent une même surface
observable : la colonne \mcptool{unit_activities.status}. C'est la
\emph{seule} information de progression disponible pour un consommateur
externe en V2 — aucun streaming de tokens, aucune interruption, aucun
\gls{sse} — et elle suffit à bâtir une UI ou un \gls{sdk}
asynchrone propre.

\begin{figure}[H]
\centering
\begin{tikzpicture}[
    node distance=6mm,
    state/.style ={rectangle, rounded corners=4pt, draw=myblue!60,
                   line width=0.6pt, fill=myblue!6,
                   minimum width=2.0cm, minimum height=0.9cm,
                   font=\footnotesize\ttfamily, align=center},
    terminal/.style ={rectangle, rounded corners=4pt, draw=gray!70,
                   line width=0.6pt, fill=gray!12,
                   minimum width=2.0cm, minimum height=0.9cm,
                   font=\footnotesize\ttfamily, align=center},
    arrow/.style ={->, thick, -{Stealth[length=2.6mm]},
                   color=black!70},
    edgelabel/.style={font=\scriptsize\itshape,
                      fill=white, inner sep=1pt, color=black!70}
]
%% --- ligne supérieure : flot principal ---
\node[state]    (d) at (0,0)     {draft};
\node[state]    (q) at (2.8,0)   {queued};
\node[state]    (s) at (5.6,0)   {searching};
\node[state]    (g) at (8.4,0)   {generating};
\node[terminal] (r) at (11.2,0)  {ready};
%% --- ligne inférieure : états d'erreur ---
\node[terminal] (f) at (8.4,-2.2)  {failed};
\node[terminal] (b) at (11.2,-2.2) {blocked};

%% --- transitions du flot nominal ---
\draw[arrow] (d) -- (q);
\draw[arrow] (q) -- (s);
\draw[arrow] (s) -- (g);
\draw[arrow] (g) -- (r);

%% --- transitions vers états d'erreur ---
\draw[arrow] (s.south) .. controls +(0,-1.0) and +(-0.8,0.7) .. (f.north west);
\draw[arrow] (g) -- (f);
\draw[arrow, dashed] (f) -- (b);

%% --- boucle de retry : passe LARGEMENT sous la ligne haute ---
\draw[arrow] (f.south) .. controls +(0,-1.4) and +(0,-1.4) .. (q.south);
\node[font=\scriptsize\itshape, color=black!75] at (5.6,-3.2) {retry manuel};
\end{tikzpicture}
\caption{Machine d'état publique d'une activité en Version~2
  (colonne \texttt{unit\_activities.status}). Les états \texttt{ready},
  \texttt{failed} et \texttt{blocked} sont terminaux (grisés). Le
  rebouclage \texttt{failed} $\to$ \texttt{queued} matérialise le
  retry manuel déclenché par l'enseignant ; la transition pointillée
  \texttt{failed} $\to$ \texttt{blocked} matérialise la propagation
  d'échec aux activités aval dans le graphe de dépendances
  \texttt{depends\_on}.}
\label{fig:v2-status}
\end{figure}

\subsection{Agent 4 : Tuteur}
\label{sec:agent4}

\paragraph{Identité.} Assistant pédagogique au service de l'enseignant. Ses
fonctions : (a)~réviser les soumissions d'étudiants et rédiger un retour
constructif, (b)~générer du matériel d'appui (plans de leçon, idées
d'ateliers), (c)~suggérer des améliorations aux activités existantes,
(d)~répondre à des questions pédagogiques.

\paragraph{Outils.} En lecture, l'agent dispose des outils nécessaires
pour ancrer ses retours dans le contenu réel. Il peut ouvrir le corps
de l'activité évaluée (\mcptool{get_activity(activity_id)}) et lire
ses voisines pour le contexte
(\mcptool{list_activities_for_unity(unity_id)}), puis solliciter la
recherche sémantique
(\mcptool{find_related_activities(syllabus_id, query)}) pour proposer
des variantes d'exercice ou du matériel d'appui. Pour personnaliser
le ton, il consulte la progression de l'étudiant
(\mcptool{get_student_progress(student_id, syllabus_id?)}) ; pour
récupérer le travail à évaluer, il appelle
\mcptool{get_submissions(student_id, activity_id?)}. Côté écriture,
l'agent ne dispose que d'un seul outil :
\mcptool{save_feedback(submission_id, feedback, score)}, qui persiste
le retour rédigé à destination de l'étudiant.

\paragraph{Workflow (révision de soumission).}
La révision d'une soumission suit un déroulé stable. L'agent
commence par récupérer le travail de l'étudiant via
\mcptool{get_submissions(student_id, activity_id)}, puis lit les
objectifs de l'activité correspondante avec
\mcptool{get_activity(activity_id)} afin d'aligner son retour sur ce
qui était réellement attendu. De façon optionnelle, il consulte la
progression générale de l'étudiant
(\mcptool{get_student_progress(student_id)}) pour ajuster son ton
au niveau de l'apprenant. Vient ensuite la rédaction proprement
dite, qui s'organise en trois temps : identification d'un à trois
points forts spécifiques (citant la réponse réelle de l'étudiant),
identification d'un à trois axes d'amélioration tout aussi
spécifiques, puis suggestion d'une prochaine étape concrète —
activité à retravailler, exercice voisin à tenter, concept à
relire. L'agent termine par un appel à
\mcptool{save_feedback(submission_id, feedback, score)}, qui
persiste l'ensemble dans la base SprintUp.

\paragraph{Schéma de sortie.} Pour une révision : objet \gls{json} avec
\mcptool{submission_id}, \mcptool{score} (0--100), \mcptool{strengths[]},
\mcptool{improvements[]}, \mcptool{next_step},
\mcptool{feedback_persisted: true}. Pour un plan de leçon : objet \gls{json}
avec \mcptool{topic}, \mcptool{grade_level}, \mcptool{duration_min},
\mcptool{objectives[]}, \mcptool{sections[]} (titre, durée, activités),
\mcptool{materials[]}.

\paragraph{Barre de qualité.} Trois exigences caractérisent un
retour de qualité. Il doit d'abord être constructif et encourageant
plutôt que sévère, faute de quoi il échouera à motiver l'étudiant
à reprendre son travail. Il doit ensuite être \emph{spécifique}, en
citant la réponse réelle de l'étudiant plutôt qu'en se contentant
de formules génériques applicables à n'importe quelle copie. Pour
les plans de leçon, le cumul des durées des sections doit enfin
totaliser exactement la durée demandée dans la requête.

\paragraph{Règles.} Trois interdits encadrent l'agent. Il ne doit
\emph{jamais} rédiger un retour sans avoir au préalable lu les
objectifs de l'activité concernée, sous peine de tomber dans le
commentaire décontextualisé. Il ne doit \emph{jamais} attribuer un
score en dehors de la plage 0--100. Il ne doit enfin \emph{jamais}
exposer les données d'un autre étudiant : le retour est strictement
attaché au couple (student\_id, submission\_id) reçu en entrée.

\subsection{Agent 5 : Assistant élève}
\label{sec:agent5}

\paragraph{Identité.} Agent d'accompagnement au service de l'apprenant. Son
rôle est d'aider l'étudiant à apprendre, en expliquant des concepts, en
échafaudant sa réflexion et en recommandant la suite du parcours. Ton
chaleureux, patient et encourageant.

\paragraph{Outils (lecture seule).} L'assistant élève opère
strictement en lecture sur le serveur MCP : il ne dispose d'aucun
outil d'écriture, et n'a donc aucun moyen de modifier la base de
données de la plateforme. Sa trousse d'outils couvre essentiellement
la situation de l'étudiant
(\mcptool{get_student_progress(student_id, syllabus_id?)}), le
curriculum complet (\mcptool{get_syllabus(syllabus_id)}) et son
balayage (\mcptool{list_unities(syllabus_id)},
\mcptool{list_activities_for_unity(unity_id)}). Lorsque l'étudiant
bloque sur un concept, l'agent ouvre directement le corps de
l'activité concernée (\mcptool{get_activity(activity_id)}), et,
le cas échéant, mobilise la recherche sémantique
(\mcptool{find_related_activities(syllabus_id, query)}) pour
proposer un exemple alternatif ou une activité voisine éclairante.

\paragraph{Workflow.}
L'agent commence par identifier l'intention de l'étudiant — demande
d'explication, d'indice, de résumé ou recommandation «~que faire
ensuite~?~». Il récupère ensuite la progression de l'étudiant pour
personnaliser sa réponse au niveau réel de l'apprenant. Si
l'étudiant pose une question sur un exercice ou un quiz, l'agent
ne révèle \emph{jamais} la réponse directe : il reformule le
problème en termes plus simples, pointe vers le concept ou
l'exemple résolu pertinent, et pose une question socratique pour
guider l'étudiant vers l'étape suivante. Lorsque la demande est
plutôt «~que dois-je faire ensuite~?~», l'agent lit la progression
de l'étudiant, identifie la prochaine activité non terminée et la
recommande explicitement.

\paragraph{Schéma de sortie.} Réponse conversationnelle en prose (markdown
autorisé). Pour une recommandation «~que faire ensuite~», ajouter un bloc
\gls{json} :
\begin{center}
\texttt{\{"recommendation":\,\{"activity\_id":\,"...",\,"title":\,"...",\,"why":\,"..."\}\}}
\end{center}

\paragraph{Barre de qualité.} Trois exigences gouvernent la qualité
des échanges. Le ton doit rester chaleureux et patient, sans jamais
basculer dans la condescendance. Les explications doivent être
adaptées au niveau effectif de l'étudiant, ce que la lecture de sa
progression permet de calibrer. Enfin, lorsque l'agent fournit un
indice, celui-ci doit \emph{conduire} à la réponse sans la révéler :
c'est la condition nécessaire pour que l'échange reste pédagogique
et non simplement transactionnel.

\paragraph{Règles.} L'assistant obéit à trois interdits stricts. Il
ne donne \emph{jamais} une réponse directe à un exercice, une
question ou un quiz, même face à une demande explicite — c'est ce
qui le distingue d'un simple chatbot généraliste. Il ne révèle
\emph{jamais} les données d'un autre étudiant que celui qui
l'interroge. Il ne recommande \emph{jamais} une activité en dehors
du syllabus courant de l'étudiant, ce qui protège la cohérence du
parcours défini par l'enseignant SprintUp.

\subsection{Extension envisagée : assistant administrateur}
\label{sec:agent6}

L'assistant administrateur n'est pas inclus dans la première version livrée.
Sa conception, esquissée ci-dessous, reprend les principes des agents en
lecture seule et sera détaillée dans une version ultérieure.

\paragraph{Identité.} Agent d'aide à la décision pour les administrateurs
de la plateforme. Génère des tableaux de bord, des rapports agrégés et des
alertes sur l'état du contenu et de l'engagement.

\paragraph{Outils (lecture seule, serveur Analytique).}
\mcptool{get_enrollment_report},
\mcptool{get_completion_report},
\mcptool{get_content_report},
\mcptool{get_user_stats},
\mcptool{list_students},
\mcptool{list_syllabuses}.

\paragraph{Règles clés.} Deux interdits sont déjà acquis par
conception. L'agent ne doit jamais exposer de données personnelles
d'étudiants sans autorisation explicite, conformément au principe
du moindre privilège appliqué au serveur Analytique. Il ne doit
par ailleurs jamais inventer de métrique : lorsqu'un outil retourne
une valeur nulle ou indisponible, il en signale explicitement la
nature plutôt que de fabriquer un chiffre intermédiaire.

%% ────────────────────────────────────────────────────────────
\section{Lecture comparative des implémentations}
\label{sec:comparatif}

Cette section ne cherche pas à désigner une variante gagnante : elle
relit les implémentations livrées comme autant de \emph{réponses
différentes à des questions d'intégration différentes}. Elle montre
ce que chaque architecture privilégie (contrôle opérateur,
auditabilité, réactivité, intégrabilité programmatique) et ce qu'elle
paie en contrepartie. Les observations chiffrées rapportées au fil du
texte ont été recueillies sur l'instance déployée à des fins
d'\emph{ordres de grandeur} ; elles éclairent les contraintes
pratiques de chaque variante, sans constituer un test de conformité.

\subsection{Axes de comparaison}
\label{sec:axes-compar}

Quatre axes structurent la lecture. \textbf{Propriétaire de
l'orchestration} : qui décide, à l'exécution, du prochain pas du
graphe — un superviseur \gls{llm} ou du code \textsc{TypeScript}
déterministe ? Cet axe gouverne directement l'auditabilité et la
reproductibilité. \textbf{Posture humaine} : la variante suppose-t-elle
un enseignant disponible en temps réel (\emph{human-in-the-loop}),
un opérateur qui pilote la production de manière asynchrone via une
machine d'état, ou aucun humain (déclencheur programmatique) ?
\textbf{Granularité de l'artefact inspectable} : le travail
intermédiaire transite-t-il par un flux de tokens éphémère, par des
fichiers persistants dans un \gls{vfs} partagé, ou par des lignes
typées en base de données ? \textbf{Coût et latence} : l'enveloppe
de tokens et le temps de bout en bout situent la variante par
rapport à l'attente de l'utilisateur final, plus qu'ils ne la
qualifient en valeur absolue.

\subsection{Agent~1 contre agent~2 : la valeur conditionnelle de l'ancrage}
\label{sec:a1-vs-a2}

Le générateur d'activités sans \gls{mcp} (agent~1) et son homologue
outillé (agent~2) partagent un même schéma de sortie mais diffèrent
radicalement par leur posture vis-à-vis de l'existant. Sur les sujets
rejoués depuis l'instance déployée, deux régimes se distinguent.
Lorsque le sujet recoupe le syllabus lié, l'agent~2 mobilise ses
outils de lecture, consulte les activités sœurs et produit un contenu
\emph{ancré} dans les objectifs et le vocabulaire de l'unité
concernée ; l'agent~1, par construction, ne peut produire qu'un
contenu générique. Lorsque le sujet sort du périmètre, l'agent~2
ne déclenche aucune lecture et se replie sur le comportement de
l'agent~1. Ce repli est un comportement de conception, non un
défaut~: il signifie que le coût d'ancrage n'est payé que quand il
apporte une valeur.

L'ancrage a précisément un coût observable. Sur les sujets rejoués
depuis l'instance déployée, et à titre d'ordre de grandeur,
l'agent~2 consomme environ trois fois plus de
tokens d'entrée que l'agent~1 et présente un surcoût de latence de
l'ordre de quelques dizaines de pour cent, entièrement imputable à
la phase de lecture qui précède la génération. Les tokens de sortie
restent comparables : le surcoût n'est pas dans la rédaction, mais
dans le contexte qu'elle exige. Cet écart n'est pas un défaut à
corriger~; il est le tarif explicite de l'ancrage. Il pose, en
revanche, la question architecturale du \emph{moment} où ce coût
doit être payé — à chaque génération, ou une fois par session — que
le chapitre suivant reprend dans sa partie consacrée au cache de
contexte.

L'enseignement principal de cette paire est qualitatif :
l'indexation et le \gls{mcp} n'apportent pas une amélioration
uniforme mais une amélioration \emph{conditionnelle au sujet}. Une
stratégie d'orchestration sensible à ce conditionnement (router vers
l'agent~1 ou vers l'agent~2 selon la pertinence du périmètre)
préserverait la qualité d'ancrage sans en payer le coût quand il
n'apporte rien : c'est l'un des leviers que le chapitre~3 explore.

\subsection{Quatre variantes de l'agent~3, quatre profils d'intégration}
\label{sec:4variantes}

Les quatre variantes de l'agent~3, présentées à la
\textsection\ref{sec:agent3}, se distinguent moins par la qualité
intrinsèque du contenu généré que par la \emph{posture} qu'elles
supposent côté opérateur et côté intégrateur. Nous les relisons ici
à travers les axes définis plus haut.

\paragraph{V1/Op.~1 — conversation + \emph{Human-in-the-Loop}.}
Cette variante place le superviseur \gls{llm} au cœur du routage :
c'est le modèle qui décide, à chaque tour, d'enchaîner intake,
recherche, écriture ou réplique. Le contrat avec l'enseignant est
celui d'un dialogue : interruptions structurées (formulaire d'intake,
questions de clarification), streaming de tokens, possibilité de
corriger en ligne. Cette approche est mieux adaptée à un enseignant
qui souhaite \emph{co-construire} un syllabus à voix haute, ou
clarifier interactivement une commande mal spécifiée. En contrepartie,
l'orchestration n'est auditable qu'a posteriori, par les traces
\gls{llm} : aucun pipeline de code n'est exécuté de bout en bout, ce
qui limite la reproductibilité et complique l'intégration dans une
chaîne CI/CD ou un workflow industriel.

\paragraph{V1/Op.~2 — sous-agents spécialistes + VFS.}
La variante \mcptool{deepagent} conserve le superviseur \gls{llm}
mais externalise les phases en sous-agents porteurs d'instructions
système dédiées (planner, writer, activity-maker, critic). Le canal
d'inter-communication, un \gls{vfs} partagé, matérialise des
artefacts intermédiaires inspectables (\mcptool{/pedagogy\_plan.md},
\mcptool{/lessons/<id>.md}, \mcptool{/critiques/<target>.md}). Cette
décomposition est mieux adaptée à un usage où l'on souhaite
\emph{séparer explicitement les responsabilités} de chaque phase
et inspecter \emph{a posteriori} les artefacts produits par chacune.
En contrepartie, l'orchestration de plus haut niveau reste portée
par le superviseur \gls{llm}, et chaque délégation paie l'isolement
de contexte par un briefing exhaustif à reconstruire : le coût en
tokens y est sensiblement plus élevé qu'en V1/Op.~1.

\paragraph{V2/Op.~1 — atelier formulaire + ordonnanceur déterministe.}
La variante \mcptool{unit-bulk-supervisor} déplace l'orchestration
hors du \gls{llm} : aucun superviseur \gls{llm} n'intervient au
niveau du syllabus, et l'ordonnancement par vagues de dépendances est
porté par du code \textsc{TypeScript} pur. La progression de chaque
activité est exposée par une machine d'état publique
(figure~\ref{fig:v2-status}), ce qui rend la surface consommable
sans interpréter un protocole de streaming. Cette variante est mieux
adaptée à un \emph{atelier rédactionnel asynchrone}, où l'opérateur
décide quand lancer une vague et observe la progression à travers
des statuts typés. En contrepartie, l'absence d'interruption en
ligne et de streaming impose une interface plus dense côté
enseignant, et exclut par construction la correction conversationnelle
en cours de génération.

\paragraph{V2/Op.~2 — orchestrateur SDK \emph{one-shot}.}
La variante \mcptool{sdk-orchestrator} pousse le déplacement plus
loin : l'unique surface est un endpoint REST qui répond
\textsc{HTTP}~202 et exécute la génération en arrière-plan, en
parallélisant les unités via la primitive \mcptool{Send} de
LangGraph. La planification produit un contrat de cours complet
\emph{avant} la première écriture, ce qui réduit la dérive
structurelle. Cette variante est mieux adaptée à un
\emph{déclenchement industriel} (autre application, batch nocturne,
ingestion programmatique). En contrepartie, l'enseignant n'a aucun
levier pendant l'exécution : la seule observation est la machine
d'état publique, comme en V2/Op.~1, et toute correction suppose de
repasser par la surface atelier.

\subsection{Coûts observés et enveloppes d'usage}
\label{sec:observations}

Quelques mesures issues de l'instance déployée fixent les ordres de
grandeur, sans constituer un classement. La génération d'un syllabus
complet via la V1/Op.~1 prend, selon la complexité de la requête et
l'activation ou non d'un cycle critique, de l'ordre de la
demi-minute à deux minutes. La production d'une activité par
l'agent~2 paie une phase de lecture qui ajoute de l'ordre de
plusieurs dizaines de pour cent de latence à la ligne de base de
l'agent~1, et environ triple la taille du prompt d'entrée. Les voies
V2, qui n'exposent pas de latence interactive, se caractérisent
plutôt par leur \emph{débit agrégé} d'une vague, lequel dépend du
parallélisme inter-unités et du plafond imposé par l'ordonnanceur.

Ces chiffres ne désignent pas une variante préférable : ils
qualifient des \emph{enveloppes d'usage} différentes. Une enveloppe
interactive (V1) suppose une attente d'utilisateur de l'ordre de la
minute et tolère le streaming. Une enveloppe atelier (V2/Op.~1)
suppose une production asynchrone observable par statut. Une
enveloppe industrielle (V2/Op.~2) suppose un déclencheur
programmatique non bloquant. Chaque enveloppe est cohérente avec
sa variante d'origine et incohérente avec les autres : c'est ce
constat, plus que les valeurs absolues mesurées, qui motivera le
chapitre suivant.

\subsection{Enseignements croisés}
\label{sec:enseignements}

Trois enseignements émergent de cette lecture comparative. D'abord,
la \emph{propriété de l'orchestration} (\gls{llm} contre code) est
l'axe structurant : il décide de l'auditabilité, de la
reproductibilité et du coût d'intégration plus que le nombre de
sous-agents ou la nature du critic. Ensuite, les variantes diffèrent
moins par la qualité du contenu que par la \emph{posture humaine}
qu'elles supposent ; aucune n'est universellement meilleure parce
qu'aucune ne s'adresse à la même phase du parcours pédagogique de
SprintUp. Enfin, la valeur de l'indexation et du \gls{mcp} n'est pas
uniforme : elle est conditionnelle au périmètre de la requête, ce
qui suggère qu'elle relève d'une stratégie de \emph{routage} et non
d'une activation inconditionnelle.

Ces trois constats — propriété de l'orchestration, posture humaine,
conditionnalité de l'ancrage — ne se résolvent pas au niveau d'une
variante isolée. Ils appellent une \emph{stratégie d'orchestration}
qui articule la coexistence des variantes selon la phase du
workflow et le contexte d'intégration. Le chapitre suivant prendra
en charge cette stratégie : routage entre variantes, articulation des
surfaces conversationnelle, atelier et SDK, politique de sécurité
adossée au principe du moindre privilège, et procédure
d'opérationnalisation sur l'instance déployée.

%% ────────────────────────────────────────────────────────────
\section{Conclusion : du catalogue à l'orchestration}
\label{sec:ch2-conclusion}

Ce chapitre a présenté la conception logique de la couche d'agents
de SprintUp et la lecture comparative des implémentations livrées.
Trois éléments le structurent. \textbf{L'indexation vectorielle} ne
remplace pas la lecture directe des activités sœurs : elle la
prolonge, comme garde-fou intra-syllabus et comme instrument
principal d'anti-redondance et de recommandation à l'échelle de
toute la base. \textbf{Le protocole \gls{mcp}} standardise l'accès
aux données et matérialise concrètement le principe du moindre
privilège : trois serveurs disjoints, une liste blanche d'outils par
agent, aucun chemin direct de l'agent vers la base de données.
\textbf{Le catalogue des agents} — deux générateurs d'activités,
quatre variantes du concepteur de syllabus, un tuteur, un assistant
élève et un assistant administrateur — couvre l'ensemble du cycle
pédagogique sans imposer un unique paradigme d'interaction.

La lecture comparative de la \textsection\ref{sec:comparatif} a mis
en évidence le point central de ce chapitre : les quatre variantes
de l'agent~3 ne s'inscrivent pas sur une échelle linéaire de
qualité, mais répondent à des questions d'intégration distinctes
(co-construction conversationnelle, décomposition spécialiste,
atelier asynchrone, déclenchement industriel). Aucune n'est
universellement préférable, parce qu'aucune ne s'adresse à la même
phase du workflow ni au même contexte d'usage. Cette pluralité est
un constat de conception, pas une indécision : elle reflète la
diversité réelle des cas d'usage que SprintUp doit servir.

Cette pluralité ouvre, en retour, une question d'\emph{orchestration
de niveau supérieur}. Comment router une demande pédagogique vers la
variante adaptée à la phase du parcours et au contexte d'intégration ?
Comment articuler les surfaces conversationnelle, atelier et SDK
sans démultiplier le code de plomberie ? Comment opérationnaliser,
observer et sécuriser la coexistence de ces variantes sur l'instance
déployée ? Et comment, sur cette même couche, traduire le principe
du moindre privilège du \gls{mcp} en une politique de sécurité
applicative au sens propre ?

Le chapitre~3 prendra ces questions en charge. Il définira la
stratégie d'orchestration des variantes (politique de routage,
articulation des surfaces, mutualisation des sous-graphes communs),
décrira le dispositif d'observabilité associé (traces \gls{llm},
machine d'état publique, métriques de débit et de latence), et y
adossera l'analyse de sécurité — surface d'attaque, contrôle
d'accès, gestion des secrets, traitement des données personnelles —
ainsi que les procédures d'opérationnalisation nécessaires à la
mise en production. Le chapitre~2 a comparé les formes d'agents ; le
chapitre~3 en organisera la coexistence.

\newpage
